\documentclass[11pt]{article}
\usepackage[margin=1in]{geometry}
\usepackage{amsmath,amssymb,amsthm,hyperref}
\newtheorem{theorem}{Theorem}
\newtheorem{lemma}{Lemma}
\newtheorem{proposition}{Proposition}
\newcommand{\HH}{\mathsf H}
\newcommand{\II}{\mathsf I}
\newcommand{\E}{\mathbb E}
\newcommand{\R}{\mathbb R}
\title{Directional omission weights:\\an exact feasibility criterion and a nonregular entropy extension}
\author{Current-literature subagent: independent derivation}
\date{November 15, 2026}
\begin{document}
\maketitle

\paragraph{Scope.}
This note proves a sufficient entropy mechanism for nonregular fixed
degree frames. It does not assert recovery for every minimum-degree-three
base. In particular the $K_{3,12}$ information obstruction lies outside
the feasible-weight class. The feasibility result is a capacitated
fractional-factor calculation, proved here without a priority claim.
Non-strict feasibility gives the first-moment bound. The qualitative
one-hot rigidity assertion below uses the explicitly stronger strict
feasibility assumption.

\paragraph{Frame.}
Fix a nonnegative integer matrix $Q=(d_{ij})_{i,j\in[k]}$ and positive
weights $\alpha_i$ with $\sum_i\alpha_i=1$ and
$\alpha_i d_{ij}=\alpha_jd_{ji}$. Write $d_i=\sum_jd_{ij}$.
An internal factor $d_{ii}$ counts physical stubs, so an internal
undirected edge has half the corresponding directed mass. In the
configuration model, $n_i=\alpha_i n$ are integers,
$n_id_{ii}$ are even, and all active internal pairings and cross
bijections are independent and uniform. A loop contributes two to
degree and adjacency. Zero factors are absent.

\begin{theorem}[Exact omission-weight criterion]\label{thm:lp}
The following are equivalent.
\begin{enumerate}
\item There are nonnegative directional numbers $q_{ij}$ on active
factors such that
\[
 \sum_jd_{ij}q_{ij}=1,\qquad
 \sum_jd_{ij}q_{ji}=1,\qquad q_{ij}+q_{ji}\leq1.       \tag{1}
\]
For an internal factor the last condition is $2q_{ii}\leq1$.
\item There are symmetric numbers $q_{ij}=q_{ji}\in[0,1/2]$
with $\sum_jd_{ij}q_{ij}=1$.
\item For every $S\subseteq[k]$, define, for $i\notin S$,
\[
 d_i(S)=d_i-\sum_{j\in S}d_{ij},\qquad
 I_a(S)=\{i\notin S:d_i(S)=a\}\quad(a=0,1).
\]
Then
\[
 2\alpha(I_0(S))+\alpha(I_1(S))\leq2\alpha(S).       \tag{2}
\]
\end{enumerate}
A violating set $S$ in (2) is an explicit infeasibility certificate.
For a simple equal-fiber base $H$, (2) becomes
\[
 2\,i(H-S)+n_1(H-S)\leq2|S|\quad\text{for every }S, \tag{3}
\]
where $i$ and $n_1$ count isolated vertices and degree-one vertices.
In this case the equivalent edge variables $y_e=2q_e$ form a
capacitated fractional $2$-factor:
$0\leq y_e\leq1$ and $\sum_{e\ni v}y_e=2$.
\end{theorem}
\begin{proof}
Symmetrizing $q_{ij}$ to $(q_{ij}+q_{ji})/2$ preserves each row
normalization by the two equalities in (1); it preserves the cap and
nonnegativity. Conversely symmetric feasible weights satisfy (1).

Build a type graph with one edge for every active unordered pair and
one loop for every positive $d_{ii}$. Set
\[
 u_{ij}=\alpha_id_{ij}\ (i<j),\qquad
 u_{ii}=\alpha_id_{ii}/2,\qquad b_i=2\alpha_i.
\]
Let $A$ be its unsigned incidence matrix, counting a loop twice.
The substitution $z_e=2u_eq_e$ identifies the symmetric system with
$Az=b$, $0\leq z\leq u$. By separation from the compact convex set
$A[0,u]$, this is equivalent to
\[
 \sum_i b_it_i\leq\sum_{e=ij}u_e(t_i+t_j)_+
                    \quad\text{for every }t\in\R^k.       \tag{4}
\]
This also supplies a general real-valued dual certificate.

It suffices to check (4) for vectors with entries in $\{-1,0,1\}$.
Indeed, put $t_i(s)=\operatorname{sign}(t_i)
\boldsymbol1_{\{|t_i|\geq s\}}$. Both $t_i=\int_0^\infty t_i(s)ds$
and
\[
 (t_i+t_j)_+=\int_0^\infty(t_i(s)+t_j(s))_+\,ds
\]
hold. The second identity follows directly for equal signs; for
opposite signs the nested threshold intervals give
$(|t_i|-|t_j|)_+$, with the positive coordinate taken first.
It also covers loops. Thus integration of the signed-indicator
inequalities gives every real-vector inequality.

Take $t=1$ on $T$, $t=-1$ on $S$, and $t=0$ otherwise, where $S,T$
are disjoint. Edges inside $T$ contribute twice, edges from $T$ to
the remaining zero set once, and edges from $T$ to $S$ zero.
Consequently (4) is exactly
\[
 2\alpha(S)+\sum_{i\in T}\alpha_i(d_i(S)-2)\geq0.    \tag{5}
\]
For fixed $S$, the minimum over $T\subseteq[k]\setminus S$ is attained
by all degree-zero and degree-one indices; degree-two terms vanish.
This proves (2), both necessarily and sufficiently. Equal $\alpha_i$
give (3).
\end{proof}

\paragraph{Why an ordinary fractional perfect matching is insufficient.}
For a simple base, $q$ has weighted degree one but also the edge cap
$q_e\leq1/2$. This extra cap is substantive. Take two disjoint
$K_{3,4}$ and add a bridge between vertices on their size-four sides.
The resulting connected graph has minimum degree three and an integral
perfect matching: use the bridge and match the residual $K_{3,3}$ in
each copy. Every fractional perfect matching must put mass one on the
bridge, by subtracting the vertex demands on the two sides of either
copy. It therefore fails the omission-weight cap. For the size-three
side $S$ of one copy, $H-S$ has three isolated vertices and one
degree-one vertex, giving $2i+n_1=7>6=2|S|$.
For $K_{r,b}$ with $r<b$, taking the size-$r$ side gives
$2b>2r$ directly. Equal-fiber unbalanced bipartite bases fail (1).

\paragraph{Strict feasibility.}
Strictness means $q_{ij}>0$ on every active directed factor and
$q_{ij}+q_{ji}<1$ on every active unordered factor, including
$2q_{ii}<1$. Symmetrization preserves strictness. It can be tested by
maximizing $\theta$ subject to symmetric weights,
\[
 \sum_jd_{ij}q_{ij}=1,\qquad
 \theta\leq q_{ij}\leq(1-\theta)/2.                 \tag{6}
\]
The optimal $\theta$ is positive exactly under strict feasibility.
Non-strict feasibility in Theorem~\ref{thm:lp} alone does not imply
this. Regular frames with $d\geq3$ have $q_{ij}=1/d$ and are strict.
There are also strictly feasible nonregular simple bases. In
$K_5$ minus one edge, let $A$ be the two endpoints of the deleted
edge and $B$ the other three vertices. Give every $A$--$B$ edge
weight $1/3$, and every $B$--$B$ edge weight $1/6$. Each $A$ vertex
has weight sum $3(1/3)=1$, and each $B$ vertex has sum
$2(1/3)+2(1/6)=1$. Every edge weight is positive and below $1/2$;
the vertex degrees are three and four.

\paragraph{Entropy setup.}
Use any finite abelian vector alphabet $\mathcal A$, with a fixed local
map $f:\mathcal A\to\mathcal A$. At an old type-$i$ vertex an ordered
star $(X_i,Z_{i1},\ldots,Z_{id_i})$ is legal when
\[
 \sum_{s=1}^{d_i} Z_{is}=f(X_i).                    \tag{7}
\]
More generally all that is needed is unique determination of any
omitted slot by the root and retained slots. The slot types are fixed
by $Q$. Compatible pooled pair laws $\nu_{ij}$ reverse on exchanging
$i,j$, with root marginals $\rho_i,\rho_j$. Define
\[
 w_{ii}=\alpha_id_{ii}/2,\qquad w_{ij}=\alpha_id_{ij}\ (i<j),
 \qquad
 \Phi(P)=\sum_i\alpha_i\HH(P_i)-\sum_{i\leq j}w_{ij}\HH(\nu_{ij}).
                                                               \tag{8}
\]
Absent factors are always omitted from sums.
Symmetrize the star law only over permutations of slots with the same
old neighbor type, writing $\overline P_i$. This preserves legal
support and pooled pair laws, and raises star entropy. Each physical
slot toward type $j$ then has pair law $\nu_{ij}$.
For a slot $t$ toward type $j$, put $q_{it}=q_{ij}$, and let
\[
 J_{it}=\sum_{s\ne t}\HH(Z_{is}\mid X_i)
              -\HH((Z_{is})_{s\ne t}\mid X_i)\geq0,
 \qquad I_{ij}=\II(\nu_{ij}).
\]

\begin{theorem}[Weighted all-slot identity]\label{thm:entropy}
Assume the outgoing equalities $\sum_jd_{ij}q_{ij}=1$, with
$q_{ij}\geq0$. Then the exact identity is
\begin{align}
 \Phi(\overline P)
 ={}&\sum_i\alpha_i\left(1-\sum_jd_{ij}q_{ji}\right)\HH(\rho_i)
 \nonumber\\
 &-\sum_{i\leq j}w_{ij}(1-q_{ij}-q_{ji})I_{ij}
 -\sum_i\alpha_i\sum_{t=1}^{d_i}q_{it}J_{it}.       \tag{9}
\end{align}
For an internal factor the mutual-information coefficient is
$1-2q_{ii}$. Under (1), every right-hand term is nonpositive, and
$\Phi(P)\leq\Phi(\overline P)\leq0$.
\end{theorem}
\begin{proof}
Deletion of any physical slot is lossless under (7). Average the exact
identity for deleting slot $t$ with weights $q_{it}$, whose sum is one:
\begin{align*}
 \HH(\overline P_i)
 &=\HH(\rho_i)+\sum_s(1-q_{is})\HH(Z_{is}\mid X_i)
                         -\sum_tq_{it}J_{it}\\
 &=(2-d_i)\HH(\rho_i)
       +\sum_jd_{ij}(1-q_{ij})\HH(\nu_{ij})-\sum_tq_{it}J_{it}.
\end{align*}
The coefficient $2-d_i$ uses the outgoing normalization.
Inserting into (8), reversibility makes the coefficient on each
undirected pair entropy $w_{ij}(1-q_{ij}-q_{ji})$.
Replace that entropy by
$\HH(\rho_i)+\HH(\rho_j)-I_{ij}$, with two copies of the
root entropy when $i=j$. At vertex $i$ the total root coefficient is
\[
 \alpha_i\left[(2-d_i)+\sum_jd_{ij}(1-q_{ij}-q_{ji})\right]
 =\alpha_i\left[1-\sum_jd_{ij}q_{ji}\right].
\]
This proves (9). Entropy concavity proves the unsymmetrized inequality.
\end{proof}

\begin{proposition}[Explicit nonregular modular moment]\label{prop:moment}
Suppose (1) is feasible, $k,r\leq K$, and $\max_i d_i\leq D$.
Let $p$ be prime, $R\in\mathbb F_p^{r\times r}$ arbitrary, and
$Z_R(A)=|\{Y\in\mathbb F_p^{n\times r}:AY=YR\}|$.
For every admissible configuration size, set
\[
 q_*=p^K,\qquad F_*=K+\frac{K(K+3)}2q_*^2,
 \qquad B_*=Kq_*^D+F_*.
\]
Then
\[
 \E Z_R(A)\leq(3D)^{F_*}(n+1)^{B_*}.              \tag{10}
\]
For each fixed frame, conditioning all factors to be simple costs a
constant factor along its admissible size sequence. No lower bound on
internal degrees is required.
\end{proposition}
\begin{proof}
The exact compatible-star factorial count in the regular audit applies
word for word with the local degree $d_i$ in place of $d$: the vertex
factor is $n_i!/\prod_sT_i(s)!$, internal categories contribute
\[
 \frac{\prod_{x<y}m_{ii}(x,y)!\prod_x(m_{ii}(x,x)-1)!!}
 {(n_id_{ii}-1)!!},
\]
and cross factors contribute
$\prod_{x,y}m_{ij}(x,y)!/(n_id_{ij})!$.
The root assignment is already in the star multinomial.
Loops use the same two-stub convention as (7).
There are $(p^r)^{d_i}\leq q_*^D$ legal stars at old type $i$.
All counts are at most $Dn$, and the same uniform factorial estimates
give the contribution of a type at most
\[
 \exp\{n\Phi(P)+F_*[\log(Dn+1)+1]\}.
\]
Apply Theorem~\ref{thm:entropy}, then sum over at most
$(n+1)^{Kq_*^D}$ types. The inequalities $Dn+1\leq D(n+1)$ and
$e<3$ give (10). For fixed bounded degrees and positive proportions,
the finitely many independent regular and biregular factors each have
a positive limiting simplicity probability, including probability one
for absent or matching factors. Dividing by their product proves the
conditioning claim.
\end{proof}

\begin{proposition}[Qualitative one-hot rigidity under strict weights]
Assume strict feasibility and $d_i\geq3$ for every $i$. For real
one-hot labels $X_i=e_a$, impose the integer local equation
$\sum_s Z_{is}=X_iR$ with fixed $R$. Every compatible collection of
star laws with $\Phi(P)=0$ has every root marginal $\rho_i$ a point
mass. No distinct-row or nonsingularity assumption on $R$ is needed.
\end{proposition}
\begin{proof}
After symmetrization, nonpositivity and $\Phi(P)=0$ force
$I_{ij}=0$ on every active edge and $J_{it}=0$ for every slot.
Fix a type $i$, write $\mu_s=\E Z_{is}$ and
$V_s=\E\|Z_{is}-\mu_s\|_2^2$. The pair independence implies that
each slot marginal is independent of the root. Vanishing $J_{it}$
implies conditional independence of all retained slots given the root.
Consequently their joint law equals the product of their marginals,
independently of the root. Under this actual product law the omitted
label is $X_iR-\sum_{s\ne t}Z_{is}$. Compute
\[
 V_t=\sum_{s\ne t}V_s+B_i,
 \qquad B_i=\E\left\|X_iR-\sum_s\mu_s\right\|_2^2\geq0.
\]
The left side uses the actual omitted marginal; the right side follows
by conditioning on the root and independence of retained slots.
Summing over all $t$ gives
$(d_i-2)\sum_sV_s+d_iB_i=0$. Thus every neighbor marginal is a point
mass. Every old type is a neighbor at some active slot because its
degree is positive and the frame is reversible. Every $\rho_i$ is
therefore a point mass. Symmetrization preserves root laws.
\end{proof}

\paragraph{What does and does not follow.}
For a fixed strictly feasible frame and fixed finite candidate alphabet,
compactness supplies a strictly negative entropy bound away from all
old-type point maps. A full recovery theorem still needs the separate
small-disagreement argument and an algorithmic composition; neither is
asserted here. Boundary feasible weights may have zero mutual-information
coefficients or unused omissions, so the last proposition cannot be
inferred merely from non-strict (2).

Equation (9) also gives a limited side-information relaxation. If some
old root laws are already required to be point masses, their root
entropy coefficients need not cancel. Outgoing normalization, the edge
caps, and incoming sums at least one on all other types suffice for
nonpositivity. Degree labels alone do not generally pin root laws:
several old types may have the same degree, as in the $K_{3,12}$
obstruction.

\paragraph{Sources and provenance.}
The exact count and factorial estimates were checked against the
regular audit's \texttt{entropy-count/all-slot-intertwiner-moment.tex},
with SHA-256:
\par\noindent
\texttt{c41edb7451ebdf81417fe4d92d77d2c89e54c8317462fa37cf85db3df12447a9}.
\par\noindent
The weighted cancellation and feasibility proof above are separate
derivations. Fractional perfect $b$-matching polytopes are established
objects: Roger E. Behrend,
\href{https://arxiv.org/abs/1301.7356}{Fractional Perfect $b$-Matching
Polytopes. I: General Theory}, arXiv:1301.7356, especially Section~3.
That source treats loops as incident once, whereas our incidence matrix
counts an internal loop twice, and it does not by itself supply our
edge-cap specialization. The proof of Theorem~\ref{thm:lp} fixes both
conventions explicitly.
\end{document}
